\documentclass[10pt,a4paper]{amsart}
\usepackage[margin=19mm]{geometry}
\usepackage{amsmath,amssymb,mathtools}
\usepackage[scr=rsfs,cal=euler]{mathalfa}
\usepackage{xfrac}
\usepackage[unicode,hidelinks,bookmarks=false]{hyperref}
\newcommand{\ie}{i.e.\ }
\newcommand{\cf}{cf.\ }
\newcommand{\resp}{respectively\ }
\newcommand{\Subsection}{\subsection}
\providecommand{\nobreakdash}{\nobreak\hbox{-}}
\setlength{\emergencystretch}{2em}
\multlinegap0pt
\usepackage{bm}
\usepackage{bbm}
\usepackage{dsfont}
\usepackage{xspace}
\usepackage{cancel}
\usepackage{mathtools}
\usepackage{amsthm}
\makeatletter
\@ifundefined{theorem}{\newtheorem{theorem}{Theorem}}{}
\@ifundefined{corollary}{\newtheorem{corollary}[theorem]{Corollary}}{}
\@ifundefined{definition}{\newtheorem{definition}[theorem]{Definition}}{}
\@ifundefined{lemma}{\newtheorem{lemma}[theorem]{Lemma}}{}
\@ifundefined{proposition}{\newtheorem{proposition}[theorem]{Proposition}}{}
\@ifundefined{remark}{\newtheorem{remark}[theorem]{Remark}}{}
\makeatother
\DeclareMathOperator{\Tr}{Tr}
\title[On Toeplitz determinants with slow Fourier decay]{On Toeplitz determinants with slow Fourier decay}
\author{Nedialko Bradinoff, Maurice Duits}
\date{Source: arXiv:2608.13182v1 (CC BY 4.0)}
\begin{document}
\maketitle

\noindent\emph{Source and licence.} On Toeplitz determinants with slow Fourier decay. Version arXiv:2608.13182v1 (CC BY 4.0), distributed under the Creative Commons Attribution 4.0 International licence, as recorded in the arXiv licence metadata for this exact version and shown on the abstract page https://arxiv.org/abs/2608.13182v1. Full rights notice: licenses/arxiv-2608.13182-CC-BY-4.0.txt.\par\smallskip

\noindent\textit{Editorial scope.} Counted unit: the Remark whose source label is \texttt{None} in the article (source0.tex lines 2301--2303, source label \texttt{None}), delivered together with the article's own complete original proof of it (source0.tex lines 2305--2390, one \texttt{proof} environment of 86 lines). No mathematics is written, repaired or re-derived: statement and proof are the article's own text line for line. Required context retained uncounted: eq:condition\_on\_fk (lines 340--342); eq:condition\_on\_fk2 (lines 528--530); def:nehari\_complete (lines 533--542); thm:k-commutators (lines 1262--1266); cor:k-commutators (lines 1273--1300); thm:bound\_on\_difficult\_nested (lines 1831--1848); thm:generalthm (lines 2068--2077); thm:determinant\_expansion (lines 2114--2134); eqn:z(t)ksum (lines 2138--2140); lemma:barely\_coefficients (lines 2247--2252); prop:analyticity\_unbounded (lines 2277--2279); eq:unbounded\_case\_cumulant\_bound (lines 2293--2297). Every definition, display and label of those lines is retained unchanged. Omitted from this package and not redistributed: 1--339, 343--527, 531--532, 543--1261, 1267--1272, 1301--1830, 1849--2067, 2078--2113, 2135--2137, 2141--2246, 2253--2276, 2280--2292, 2298--2300, 2304--2304, 2391--2858. Nothing inside a retained excerpt is omitted or altered: every retained body is delivered exactly as the article's lines are written, so no omission receipt of any kind accompanies this package. Citations: the retained excerpts carry no citation command at all, so no bibliography entry of the article is transcribed and no reference is redistributed. Reference adaptation: none was needed and none was made; every reference inside the delivered unit resolves inside it, so no unresolved reference or citation is produced. Printed slips kept literal: no printed word, symbol, hypothesis or number of the counted statement or of its proof was altered. Layout and class: the article's own class is \texttt{amsart} with packages including mathtools, mathdots, hyperref, xcolor, bbm, subcaption; the wrapper is the lane's pinned \texttt{amsart} editorial wrapper at 10pt on A4, which restates the theorem-like environments the retained text needs and adds the article's own macro definitions that the retained text needs (\texttt{\textbackslash Tr}), with extra wrapper packages (bm, bbm, dsfont, xspace, cancel, mathtools). The delivered numbering is the unit's own. Assets: no retained span contains a figure environment, a graphics-inclusion command, a PDF-page inclusion, a table or a TikZ picture; no image file is copied into this package, the package declares no asset dependency and no image file is redistributed with it. Licence: version arXiv:2608.13182v1 of this article is distributed under the Creative Commons Attribution 4.0 International licence, as recorded in the arXiv licence metadata for this exact version and shown on the abstract page https://arxiv.org/abs/2608.13182v1. A full rights notice is delivered at licenses/arxiv-2608.13182-CC-BY-4.0.txt. Characters: every retained excerpt is pure ASCII, so the delivered engine, XeLaTeX with its default Latin Modern text font, reports no missing character.
\begin{equation}\label{eq:condition_on_fk}
|f_k| = \mathcal O(1/|k|), \qquad k \to \pm \infty.
\end{equation}

\begin{equation}\label{eq:condition_on_fk2}
|f_k|\leq \kappa |k|^{-1}, \quad\text{for}\quad k\not=0.
\end{equation}

\begin{definition}\label{def:nehari_complete}
We consider functions $f\in L^2(\mathbb{T})$ satisfying \eqref{eq:condition_on_fk2} for some $\kappa>0$ such that there exist constants $C>0$, $\kappa^*\geq \kappa>0$, and functions $g_+$, $g_-$ analytic inside/outside the unit circle and
\begin{enumerate} 
\item[i)] \label{cond:boundedness}
$\|f_++g_-\|_{\infty}, \|g_++f_-\|_{\infty}< C$
\item[ii)] \label{cond:fourier}
$|(g_{\pm})_k|\leq \frac{\kappa^*}{|k|}$ for $k\not=0$.
\end{enumerate}
We will refer to such symbols as \textbf{admissible}. 
\end{definition}

\begin{theorem}\label{thm:k-commutators}
There exists $Z\in\mathbb{C}^*[X_1,X_2,\ldots, X_k]$, which solves 
$$e^Z=e^{X_1}e^{X_2}\ldots e^{X_k}.$$
Furthermore $Z(t)= \log e^{tX_1}e^{tX_2}\ldots e^{tX_k}$ is a Lie element, that is $Z(t)$ is a linear sum of $X_1, X_2, \ldots, X_k$ and their nested commutators. 
\end{theorem}

\begin{corollary}\label{cor:k-commutators}
Let $Z(t)$ be as in Theorem \ref{thm:k-commutators}, then
\begin{multline}\label{eqn:commutator_expansion}
Z(t)
=\sum_{m=1}^{\infty}\frac{t^m}{m} 
%
\sum_{j=1}^{m}\frac{(-1)^{j+1}}{j}
%
\sum_{
\substack{
n_1+n_2+\ldots +n_j=m
\\n_1,n_2,\ldots,n_j\geq 1
}}
\\ 
\sum_{
\substack{
m_1^{(s)}+m_2^{(s)}+\ldots +m_k^{(s)}=n_s
\\
s=1,2,\ldots, j 
}
}
%
\frac
{[X_1^{m_1^{(1)}}\ldots X_k^{m_k^{(1)}}\ldots \ldots X_1^{m_1^{(j)}}\ldots X_k^{m_k^{(j)}}]}
{m_1^{(1)}! \ldots m_k^{(1)}!\ldots\ldots m_1^{(j)}! \ldots m_k^{(j)}!}.
\end{multline}
The sum of all coefficients in front of commutator words $[w]$ of length $m$ is bounded by $\frac{(ke)^m}{\sqrt{2\pi m^3}}$.
\end{corollary}

\begin{theorem}\label{thm:bound_on_difficult_nested}
Assume that the symbols $f^{(1)},f^{(2)},\ldots, f^{(m)}$ satisfy \eqref{eq:condition_on_fk2} for some common $\kappa>0$ and
are admissible, as in Definition \ref{def:nehari_complete} with underlying constants $C, \kappa^*>0$. 
%
%
%
Then
\begin{enumerate}
\item $\|\mathbf{K}_m\|_{op}\leq 
%
C^m 2^{m-1}(m-1)!,
$
\item $\|P_n\mathbf{K}_mQ_n\|_{2}, \|Q_n\mathbf{K}_mP_n\|_{2}\leq  6\kappa^* (2C)^{m-1} m! \sqrt{\log(1+m)}$,
\item $\Tr  P_n \mathbf{K}_m P_n \leq 
6 \kappa^{*2}  (2C)^{m-2} m^2\log(1+m)(m-2)!,
\quad\text{for}\quad m\geq 3.$
\end{enumerate}
\end{theorem}

\begin{theorem}\label{thm:generalthm}
Assume that $B:\ell_2\rightarrow\ell_2$ is an operator defined by an infinite matrix such that
\begin{enumerate}
\item $B$ is bounded
\item $\|P_nBQ_n\|_2, \|Q_nBP_n\|_2<\infty$, for all $n$.
\end{enumerate}
Then {for}  $ \|B\|_{op} \leq \frac{1}{3}$, $\log\det(P_ne^{B}P_n+Q_n)$ satisfies
$$\left|\log\det(P_ne^{B}P_n+Q_n)-\Tr (P_nBP_n)\right|\leq A \|P_nBQ_n+ Q_nBP_n\|_2^2,$$
where we can take $A=e^2\sum_{m=0}^{\infty} (e/3)^m (m+2)^{3/2}$.
\end{theorem}

\begin{theorem}\label{thm:determinant_expansion}
Suppose $f:\mathbb{T}\rightarrow\mathbb{C}$ satisfying \eqref{eq:condition_on_fk2}, $f_0=0$, and, moreover, there is $C>0$, $\|f_+\|_{\infty}, \|f_-\|_{\infty}< C$. Then
\begin{equation}
        C_2^{(n)}(f)= \sum_{k=1}^\infty \min (k,n) f_kf_{-k},
    \end{equation}
and for $|t|\leq 
\frac{\sqrt{13}-1}{6} 
\frac{1}{8Ce}$,
$\log\det T_n(e^{tf})$ is analytic in $t$ and
\begin{equation}
\log\det T_n(e^{tf})= t^2\sum_{k=1}^\infty \min (k,n) f_kf_{-k}+ \mathcal{O}(|t|^3),\quad\text{as}\quad n\rightarrow\infty,
\end{equation}
and, moreover,
$$|
\Psi_n(t)-t^2
\sum_{k=1}^{\infty} 
\min(k,n)
f_kf_{-k} 
|\leq \left(\frac{\kappa}{C}\right)^2(1+4A),$$
where $A>0$ is the constant in Theorem \ref{thm:generalthm}.
\end{theorem}

\begin{equation}\label{eqn:z(t)ksum}
Z(t)=\sum_{m=2}^{\infty}t^m\mathbf{K}_m^{(sum)},
\end{equation}

\begin{lemma}\label{lemma:barely_coefficients}
Suppose $f$ satisfies \eqref{eq:condition_on_fk}. Then, for sufficiently small $t$, we have $\exp({tf(z)})\in L_2(\mathbb T)$ and  the Fourier coefficients of $\exp({tf(z)})$ are analytic functions of $t$. Moreover, 
\begin{equation}\label{eqn:analytic_fc}
\left( \exp{tf(z)} \right)_k=\sum_{m=0}^{\infty}\frac{t^m}{m!}\sum_{k_1+\ldots+k_m=k} f_{k_1}\ldots f_{k_m}.
\end{equation}
\end{lemma}

\begin{proposition}\label{prop:analyticity_unbounded}
Suppose $f$ satisfies \eqref{eq:condition_on_fk}. Then for all $n\in\mathbb{N}$, $\det T_n(e^{tf})$ and $\log\det T_n(e^{tf})$ are well-defined analytic functions in a neighborhood of $t=0$. Moreover, there is an absolute constant $R_*(n)$, which bounds the radius of convergence from below, and it depends only on $n$ and the upper bound on the Fourier coefficients.
\end{proposition}

   \begin{equation}\label{eq:unbounded_case_cumulant_bound}
        \left|C_m^{(n)}(f)\right|\leq 
        \left(\frac{\kappa^*}{C}\right)^2 \left(\frac{1}{16}A +\frac{7}{24}
\right)\left(24Cm\right)^{m},
    \end{equation}

\begin{remark}
    We do not expect the bound \eqref{eq:unbounded_case_cumulant_bound} to be optimal. At various places, we have made rather crude estimates and it would be very interesting to see if, under this generality, the bound could be optimized to the form $\left|C_m^{(n)}(f)\right|\leq a c^m$ for some $c>0$, depending on $f$ but not $n$, as we proved for the bounded case.
\end{remark}

\begin{proof}
For $r<1$, we define the Poisson regularization of $f$ via $f_r(z)=f_+(rz)+ f_-(z/r),$ as in the proof of Lemma \ref{lemma:barely_coefficients}. Observe that under our assumptions $f_r, f_{r+}, f_{r-}$ are all bounded: since $|f_k|\leq \frac{K}{|k|}$, $|f_{rk}|\leq \frac{r^{|k|}K}{|k|}$, it follows that $$\|f_r\|_{\infty},\|f_+(rz)\|_{\infty},\|f_-({z/r})\|_{\infty}<-2K\log(1-r).$$


By Proposition \ref{prop:analyticity_unbounded} (assuming its notation), for $R<R_*(n)$, we have
\begin{multline}\label{eqn:approximation_C_m}
C_{m}^{(n)}(f)
= \frac{1}{2\pi i}\oint_{|t|=R}\log \det T_n(e^{tf})\frac{dt}{t^{m+1}}
\\=\frac{1}{2\pi i}\oint_{|t|=R}\log \det 
\lim_{r\uparrow 1} T_n(e^{tf_r})\frac{dt}{t^{m+1}}
\\=\lim_{r\uparrow 1}\frac{1}{2\pi i}\oint_{|t|=R}\log \det T_n(e^{tf_r})\frac{dt}{t^{m+1}}
=\lim_{r\uparrow 1} C_m^{n}(f_r).
\end{multline}
By Theorem \ref{thm:determinant_expansion}, since $f_r$ is bounded, there is a $R_{**}(r)>0$ such that for $R<\min(R_*(n), R_{**}(r))$,
\begin{multline}\label{eqn:cutoff_perform}
C_{m}^{(n)}(f_r)
= \frac{1}{2\pi i}\oint_{|t|=R}\log \det T_n(e^{tf_r})\frac{dt}{t^{m+1}}
\\=  \frac{1}{2\pi i}\oint_{|t|=R}\log \det P_n e^{Z_r(t)}P_n\frac{dt}{t^{m+1}}
\\=  \frac{1}{2\pi i}\oint_{|t|=R}\log \det P_n e^{Z_r^{(m)}(t)}P_n\frac{dt}{t^{m+1}}
,
\end{multline}
where
\begin{equation}
Z_r(t)=\sum_{p=2}^{\infty}t^p\mathbf{K}_p^{(sum)}(f_r), \text{ and for } m\geq 2,\quad Z_r^{(m)}(t)=\sum_{p=2}^{m}t^p\mathbf{K}_p^{(sum)}(f_r).
\end{equation}
The last equality in \eqref{eqn:cutoff_perform} is justified by the fact that powers of $t$ greater than $m$ do not affect the integral.
By Theorem \ref{thm:bound_on_difficult_nested}, the norms $\|\mathbf{K}_p^{(sum)}(f_r)\|_{op}$, $\|P_n\mathbf{K}_p^{(sum)}(f_r)Q_n\|_2,$ $\|Q_n\mathbf{K}_p^{(sum)}(f_r)P_n\|_2$ are bounded independently of $n$ and $r$, and so there is constant $R_{in}(m)$, independent of $n$ and $r$
(but depending on $m$) 
for which 
$$\sum_{p=2}^{m}R_{in}^p\|\mathbf{K}_p^{(sum)}(f_r)\|_{op}=1/3\quad\text{for}\quad 0<r<1.$$
By Theorem \ref{thm:generalthm}, for $|t|\leq R_{in}$
\begin{multline}
\log \det P_n e^{Z_r^{(m)}(t)}P_n
\\=  \Tr P_n Z_r^{(m)}(t)P_n+ \mathcal{O}\left( t^4\|P_n Z_r^{(m)}(t) Q_n+Q_nZ_r^{(m)}(t) P_n\|_2^2\right)
\\=t^2 \Tr P_n K_2^{(sum)}(f_r)P_n +\mathcal{O}(|t|^3)=t^2 \sum_{k=1}^{\infty}\min(n, k) r^{2k} f_k f_{-k}+\mathcal{O}(|t|^3),
\end{multline}
and the implicit bound for the remainder is independent of $n$ and $r$. The remainder term is analytic in $t$ for $R\leq R_{in}$ and has a zero of order at least $3$ at zero. Hence 
\begin{multline}
C_{m}^{(n)}(f_r)
=\frac{1}{2\pi i}\oint_{|t|=R}\log \det P_n e^{Z_r^{(m)}(t)}P_n\frac{dt}{t^{m+1}}
\\
=\frac{1}{2\pi i}\oint_{|t|=R_{in}}\log \det P_n e^{Z_r^{(m)}(t)}P_n\frac{dt}{t^{m+1}}
%
=\mathcal{O}(R_{in}^{m}).
\end{multline}
Thus, from \eqref{eqn:approximation_C_m} it follows that 
$$C_2^{(n)}(f)=\lim_{r\uparrow 1} \sum_{k=1}^{\infty} \min(n, k) r^{2k} f_k f_{-k}=\sum_{k=1}^{\infty} \min(k,n) f_k f_{-k},\quad \text{and}$$
$$C_m^{(n)}(f)=\mathcal{O}(R_{in}^m)=\mathcal{O}(1),\quad\text{as}\quad n\rightarrow\infty, \text{ for } m\geq 3.$$
In fact as in the proof of Theorem \ref{thm:determinant_expansion} we obtain an explicit bound; once again by Corollary \ref{cor:k-commutators} the sum of the coefficients in front of $m$-nested commutators in the expansion of $Z_r^{(m)}(t)$ is $\frac{(4e)^m} 
{\sqrt{2\pi m^3}}$. 
Hence by Theorem \ref{thm:bound_on_difficult_nested}, the operators in the expansion of $Z(t)$ (recall \eqref{eqn:z(t)ksum}) satisfy
\begin{enumerate}
\item 
$
\|\mathbf{K}_{\ell}^{(sum)}(f_r)\|_{op}\leq 
%
\frac{1}{2}\frac{(8Ce)^{\ell}(\ell-1)!}{\sqrt{2\pi \ell^3}},
$
\item $\|P_n\mathbf{K}_{\ell}^{(sum)}(f_r)Q_n\|_{2}, \|Q_n\mathbf{K}_{\ell}^{(sum)}(f_r)P_n\|_{2}
\leq  
3\frac{\kappa^*}{C} \frac{(8Ce)^{\ell}\ell!}{\sqrt{2\pi \ell^3}} \sqrt{\log(1+\ell)}
$,
\item 
$
\Tr  P_n \mathbf{K}_{\ell}^{(sum)}(f_r) P_n 
\leq 
\frac{3}{2}
\left(\frac{\kappa^*}{C}\right)^{2} 
\frac{(8Ce)^{\ell} \ell!}{\sqrt{2\pi \ell^3}}
\frac{\ell}{\ell-1}\log(1+\ell),
\quad\text{for}\quad \ell\geq 3.
$
\end{enumerate}
The additional Stirling approximation, $\frac{\ell!e^{\ell}}{\sqrt{2\pi \ell}}< \ell^\ell e ^{\frac{1}{12\ell}}$ gives us that for $t\leq \frac{1}{24C m}$,
$$\|Z_r^{(m)}(t)\|_{op}< \frac{1}{12},\quad \|P_nZ_r^{(m)}(t)Q_n\|_2, \|Q_nZ_r^{(m)}(t)P_n\|_2< \frac{1}{4}\frac{\kappa^*}{C},$$
$$\text{and}\quad\left|\Tr  P_nZ_r^{(m)}(t)P_n
- t^2\Tr  P_n\mathbf{K}_2^{(sum)}P_n\right|
\leq
\frac{7}{24}\left(\frac{\kappa^*}{C}\right)^2,
$$
and so as in the proof of Theorem \ref{thm:determinant_expansion}
we see that 
$$|C_m^{(n)}(f_r)|\leq \left(\frac{\kappa^*}{C}\right)^2 \left(\frac{1}{16}A +\frac{7}{24}
\right)\left(24Cm\right)^{m},$$
where $A$ is as in Theorem \ref{thm:generalthm}. Letting $r\uparrow 1$ we conclude the result.
\end{proof}

\end{document}
